\chapter{Curvas a partir de curvas}

En este capítulo se estudian las transformaciones que permiten obtener nuevas representaciones o formas geométricas a partir de una curva dada, analizando bajo qué condiciones se preservan propiedades fundamentales como la regularidad y la longitud.

\section{Reparametrizaciones}

Intuitivamente, una reparametrización será una variación del parámetro temporal de una curva, es decir, recorrerla a distinta velocidad sin alterar su traza, conservando así su geometría.

\begin{definicion}
    Consideramos una curva $\alpha:I \to \bb{R}^3$ y una aplicación $h:J\to I$. Podemos definir entonces una nueva curva $\beta:J \to \bb{R}^3$ como $\beta = \alpha\circ h$. Estaríamos en la siguiente situación:
    \begin{gather*}
        \xymatrix{
            J \ar[r]^h \ar@/_{5mm}/[rr]_{\beta=\alpha\circ h} & I \ar[r]^{\alpha} & \bb{R}^3
        }
    \end{gather*}
    A esta nueva curva $\beta$ la llamaremos \textbf{reparametrización} de $\alpha$ por medio de $h$. Además, aplicando la regla de la cadena tenemos que
    \begin{gather*}
        \beta'(t) = (\alpha \circ h)'(t) = h'(t)\alpha'(h(t))\ \ \ \forall t \in J
    \end{gather*}
\end{definicion}

\subsection{Conservación de la regularidad}

Si $\alpha:I\to \bb{R}^3$ es regular y queremos que $\beta$ también lo sea, tendremos que imponer $h'(t)\neq 0$ $\forall t \in J$. Tenemos entonces dos posibles situaciones:

\begin{itemize}
    \item Si $h'(t)>0\ \ \ \forall t \in J$, entonces $h:J\to I$ es un difeomorfismo creciente (preserva la orientación)
    \item Si $h'(t)<0\ \ \ \forall t \in J$, entonces $h:J\to I$ es un difeomorfismo decreciente (invierte la orientación)
\end{itemize}

\begin{observacion}
    Para que la traza sea idéntica, es decir, $tr(\beta) = tr(\alpha)$ necesitaremos imponer que $h$ sea sobreyectiva, ya que en dicho caso tendremos que $h(J)=I$ y entonces
    \begin{gather*}
        tr(\beta) = \beta(J) = (\alpha\circ h (J)) = \alpha(h(J)) = \alpha(I) = tr(\alpha)
    \end{gather*}
\end{observacion}

\subsection{Invarianza de la longitud}

Queremos ver ahora que la longitud de arco es una propiedad intrínseca de la curva y no depende de cómo se recorra. Para verlo calculamos la longitud de la curva en un compacto $[a,b]\subset I$.
\begin{gather*}
    L_a^b(\beta) = \int_a^b|\beta'(t)| dt = \int_a^b |h'(t)||\alpha'(h(t))| dt \overset{(\ast)}{=} \int_{h(a)}^{h(b)} |\alpha'(u)| du = L_{h(a)}^{h(b)}(\alpha)
\end{gather*}
donde en $(\ast)$ hemos aplicado el Teorema de Cambio de Variable ($u=h(t)$). Podemos concluir finalmente que la longitud es invariante por reparametrizaciones.

\section{Deformaciones}

A diferencia de la reparametrización, una deformación aplica una transformación geométrica a la curva, modificando posiblemente su traza.

\begin{definicion}
    Sea $\alpha: I \to \bb{R}^3$ una curva y $\phi:\bb{R}^3 \to \bb{R}^3$ una aplicación (normalmente un difeomorfismo), definimos la curva $\beta:I \to \bb{R}^3$ como $\beta = \phi \circ \alpha$. Estaremos en la siguiente situación:
    \begin{gather*}
        \xymatrix{
            I \ar[r]^{\alpha} \ar@/_{5mm}/[rr]_{\beta=\phi\circ \alpha} & \bb{R}^3 \ar[r]^{\phi} & \bb{R}^3
        }
    \end{gather*}
    A esta curva $\beta$ la llamaremos una \textbf{deformación} de $\alpha$ por medio de $\phi$. Su vector tangente viene dado por la aplicación lineal asociada (el diferencial):
    \begin{gather*}
        \beta'(t) = (\phi \circ \alpha)'(t) = d\phi_{\alpha(t)}(\alpha'(t))
    \end{gather*}
\end{definicion}

\subsection{Condición de regularidad}

Si queremos que la deformación de una curva regular siga siendo regular, necesitamos asegurar que el vector tangente no se anule. Si $\alpha$ es regular ($\alpha'(t) \neq 0$), entonces $\beta'(t)$ será distinto de cero si y solo si el diferencial $d\phi_{\alpha(t)}$ es inyectivo para todo $t\in I$, es decir, que su núcleo sea trivial. Por tanto, $\beta$ será regular si y solo si
\begin{gather*}
    \det(J\phi(\alpha(t))) \neq 0\ \ \ \forall t \in I
\end{gather*}
Podemos concluir entonces que si queremos que la regularidad de las curvas se conserve por deformaciones debemos usar difeomorfismos locales. Por ello, $\phi$ siempre será un difeomorfismo (tal y como se indica en la definición).\\

\begin{observacion}
    Hay difeomorfismos $\phi : \bb{R}^3 \to \bb{R}^3$ que conservan la ``geometría''. Estos son los ya estudiados movimientos rígidos. Estos serán difeomorfismos $M$ de la forma
    \Func{M}{\bb{R}^3}{\bb{R}^3}{x}{Mx = Ax + b}
    donde\footnote{$O(3)$ es el conjunto de matrices ortogonales de dimensión 3, es decir, matrices cuya inversa coincide con su transpuesta.} $A\in O(3)$, $b\in \bb{R}^3$.\\

    Si tomamos $\phi = M$ movimiento rígido, entonces al estudiar la traza de la curva deformada tenemos:
    \begin{gather*}
        tr(\beta) = M (tr(\alpha)) \equiv tr(\alpha)
    \end{gather*}
    es decir, la traza no es igual pero es equivalente (conserva la ``geometría''). Estudiando la longitud tenemos:
    \begin{gather*}
        L_a^b (\beta) = \int_a^b |\beta'(t)| dt \overset{(\ast)}{=} \int_a^b |A\alpha'(t)| dt = \int_a^b |\alpha'(t)| dt = L_a^b (\alpha)
    \end{gather*}
    donde en $(\ast)$ hemos aplicado que $\beta'(t)=A\alpha'(t)$. Tenemos entonces que los movimientos rígidos conservan la longitud.
\end{observacion}

\section{Longitud orientada de una curva}

\begin{definicion}
    Sea $\alpha : I \to \bb{R}^3$ una curva, $a\in I$. Definimos\footnote{La $S$ hace referencia a espacio, \textit{space} en inglés.} la siguiente aplicación
    \Func{S_a}{I}{\bb{R}}{t}{S_a(t) = \int_a^t |\alpha'(u)| du }
    Si estudiamos esta aplicación tendremos
    \begin{gather*}
        S_a(t) = 
        \left\{
            \begin{array}{r c c}
                L_a^t(\alpha) & \text{ si }& a\leq t\\\\
                -L_t^a(\alpha) & \text{ si } & t \leq a
            \end{array}
        \right.
    \end{gather*}
    A esta aplicación la llamaremos la \textbf{longitud orientada} de $\alpha$ respecto del instante $a\in I$.
\end{definicion}

\begin{propiedades}\
    \begin{itemize}
        \item $S_a$ es derivable y $S'_a(t) = |\alpha'(t)|$.
        \item Si $\alpha$ es regular, de la propiedad anterior podemos concluir que 
        \begin{enumerate}
            \item $S_a: I \to S_a(I)=J$ es un difeomorfismo entre intervalos abiertos.
            \item Sea $h=(S_a)^{-1} : J\to I$, teníamos $\alpha : I \to \bb{R}^3$ podemos reparametrizar $\alpha$\footnote{estamos usando la propia curva para reparametrizar la curva}. Definimos así $\beta = \alpha \circ h : J \to \bb{R}^3$ que es la reparametrización de $\alpha$ por medio del difeomorfismo $h = (S_a)^{-1}$. Sabemos además por la regla de la cadena que para todo $s \in J$ se tiene que
            \begin{gather*}
                S_a(h(s)) = s \Rightarrow 1 = h'(s)S_a'(h(s)) \Rightarrow 1 = h'(s)|\alpha(h(s))|
            \end{gather*}
            De donde concluimos que 
            \begin{gather*}
                h'(s) = \frac{1}{|\alpha'(h(s))|}
            \end{gather*}
            De esta forma, para $s\in J$ tenemos que
            \begin{gather*}
                \beta'(s) = (\alpha \circ h)'(s) = h'(s) \alpha'(h(s)) = \frac{\alpha'(h(s))}{|\alpha'(h(s))|}
            \end{gather*}
            Concluimos que 
            \begin{gather*}
                |\beta'(s)| = 1 \ \ \ \forall s \in J
            \end{gather*}
            Por tanto toda curva regular admite una reparametrización con celeridad constantemente $1$. Estas curvas verifican
            \begin{gather*}
                L_a^b(\beta) =  \int_a^b 1 dt = b-a
            \end{gather*}
            A este tipo de curvas se les llama \textbf{curvas parametrizadas por el arco}\footnote{lo notaremos por p.p.a. (no es el partido popular andaluz ni nada de eso).}
        \end{enumerate}
    \end{itemize}
\end{propiedades}

\begin{ejercicio}
    Si $\alpha : I \to \bb{R}^3$ es una curva p.p.a., ¿cuántas reparametrizaciones por el arco admite?
\end{ejercicio}
